% GENERATED FILE. Edit canonical lessons, then run npm run generate:books.
% canonical-source-hash: fnv1a64:d6f4dad7b6f82edd
% canonical-lessons: TE-C134-karu, TE-C134-saikilu, TE-C134-vimanam, TE-C134-dhanyam, TE-C134-godhumalu

\chapter{Car, Bicycle, Aeroplane, Grain, Wheat}
\label{ch:te-car-bicycle-aeroplane-grain-wheat}
\hlchaptermodality{134}

\begin{chapteropening}
\textbf{By the end of this chapter:} \emph{I can say a car, a bicycle, an aeroplane, grain and wheat.}

Complete the last lesson of chapter 134: I can say a car, a bicycle, an aeroplane, grain and wheat.
\end{chapteropening}

\section[\texorpdfstring{\te{కారు} (kāru)}{కారు (kāru)}]{\te{కారు} (kāru) — a car}

\label{lesson:TE-C134-karu}



\begin{quote}
Before the new one: say the Telugu for a boat, then the Telugu for a train.
\end{quote}

\subsection*{You'll want to know: \te{కారు}}
\textbf{\te{కారు}} — \emph{kāru} — \textquotedblleft{}a car\textquotedblright{}.

\begin{grammarlens}[title={What you've built}]
One more word for this chapter.
\end{grammarlens}

\subsection*{Your turn}
\begin{itemize}
\raggedright
  \item \emph{Say it:} \emph{kāru}
  \item \emph{Say it:} \emph{kāru}, once more
  \item \emph{Say it:} say \emph{kāru}
  \item \emph{Recall:} say the Telugu for a boat, then the Telugu for a train, then say \emph{kāru} again
\end{itemize}

\begin{tcolorbox}[breakable,colback=teal!4,colframe=teal!35!black,title={Before you move on}]
What does \te{కారు} mean? (\textquotedblleft{}A car\textquotedblright{}.) Say it once more.
\end{tcolorbox}
\section[\texorpdfstring{\te{సైకిలు} (sāikilu)}{సైకిలు (sāikilu)}]{\te{సైకిలు} (sāikilu) — a bicycle}

\label{lesson:TE-C134-saikilu}



\begin{quote}
Before the new one: say the Telugu for a train, then the Telugu for a car.
\end{quote}

\subsection*{You'll want to know: \te{సైకిలు}}
\textbf{\te{సైకిలు}} — \emph{sāikilu} — \textquotedblleft{}a bicycle\textquotedblright{}.

\begin{grammarlens}[title={What you've built}]
One more word for this chapter.
\end{grammarlens}

\subsection*{Your turn}
\begin{itemize}
\raggedright
  \item \emph{Say it:} \emph{sāikilu}
  \item \emph{Say it:} \emph{sāikilu}, once more
  \item \emph{Say it:} say \emph{sāikilu}
  \item \emph{Recall:} say the Telugu for a train, then the Telugu for a car, then say \emph{sāikilu} again
\end{itemize}

\begin{tcolorbox}[breakable,colback=teal!4,colframe=teal!35!black,title={Before you move on}]
What does \te{సైకిలు} mean? (\textquotedblleft{}A bicycle\textquotedblright{}.) Say it once more.
\end{tcolorbox}
\section[\texorpdfstring{\te{విమానం} (vimānaṁ)}{విమానం (vimānaṁ)}]{\te{విమానం} (vimānaṁ) — an aeroplane}

\label{lesson:TE-C134-vimanam}



\begin{quote}
Before the new one: say the Telugu for a car, then the Telugu for a bicycle.
\end{quote}

\subsection*{You'll want to know: \te{విమానం}}
\textbf{\te{విమానం}} — \emph{vimānaṁ} — \textquotedblleft{}an aeroplane\textquotedblright{}.

\begin{grammarlens}[title={What you've built}]
One more word for this chapter.
\end{grammarlens}

\subsection*{Your turn}
\begin{itemize}
\raggedright
  \item \emph{Say it:} \emph{vimānaṁ}
  \item \emph{Say it:} \emph{vimānaṁ}, once more
  \item \emph{Say it:} say \emph{vimānaṁ}
  \item \emph{Recall:} say the Telugu for a car, then the Telugu for a bicycle, then say \emph{vimānaṁ} again
\end{itemize}

\begin{tcolorbox}[breakable,colback=teal!4,colframe=teal!35!black,title={Before you move on}]
What does \te{విమానం} mean? (\textquotedblleft{}An aeroplane\textquotedblright{}.) Say it once more.
\end{tcolorbox}
\section[\texorpdfstring{\te{ధాన్యం} (dhānyaṁ)}{ధాన్యం (dhānyaṁ)}]{\te{ధాన్యం} (dhānyaṁ) — grain}

\label{lesson:TE-C134-dhanyam}



\begin{quote}
Before the new one: say the Telugu for a bicycle, then the Telugu for an aeroplane.
\end{quote}

\subsection*{You'll want to know: \te{ధాన్యం}}
\textbf{\te{ధాన్యం}} — \emph{dhānyaṁ} — \textquotedblleft{}grain\textquotedblright{}.

\begin{grammarlens}[title={What you've built}]
One more word for this chapter.
\end{grammarlens}

\subsection*{Your turn}
\begin{itemize}
\raggedright
  \item \emph{Say it:} \emph{dhānyaṁ}
  \item \emph{Say it:} \emph{dhānyaṁ}, once more
  \item \emph{Say it:} say \emph{dhānyaṁ}
  \item \emph{Recall:} say the Telugu for a bicycle, then the Telugu for an aeroplane, then say \emph{dhānyaṁ} again
\end{itemize}

\begin{tcolorbox}[breakable,colback=teal!4,colframe=teal!35!black,title={Before you move on}]
What does \te{ధాన్యం} mean? (\textquotedblleft{}Grain\textquotedblright{}.) Say it once more.
\end{tcolorbox}
\section[\texorpdfstring{\te{గోధుమలు} (gōdhumalu)}{గోధుమలు (gōdhumalu)}]{\te{గోధుమలు} (gōdhumalu) — wheat}

\label{lesson:TE-C134-godhumalu}



\begin{quote}
Before the new one: say the Telugu for an aeroplane, then the Telugu for grain.
\end{quote}

\subsection*{You'll want to know: \te{గోధుమలు}}
\textbf{\te{గోధుమలు}} — \emph{gōdhumalu} — \textquotedblleft{}wheat\textquotedblright{}.

\begin{grammarlens}[title={What you've built}]
One more word for this chapter.
\end{grammarlens}

\subsection*{Your turn}
\begin{itemize}
\raggedright
  \item \emph{Say it:} \emph{gōdhumalu}
  \item \emph{Say it:} \emph{gōdhumalu}, once more
  \item \emph{Say it:} say \emph{gōdhumalu}
  \item \emph{Recall:} say the Telugu for an aeroplane, then the Telugu for grain, then say \emph{gōdhumalu} again
\end{itemize}

\begin{tcolorbox}[breakable,colback=teal!4,colframe=teal!35!black,title={Before you move on}]
What does \te{గోధుమలు} mean? (\textquotedblleft{}Wheat\textquotedblright{}.) Say it once more.
\end{tcolorbox}
